\chapter{Tester et déployer des systèmes à faible latence}\label{ch:ll:testing-and-deploying-low-latency-systems}

Une version d'un système de trading peut passer tous les tests fonctionnels et rester pourtant moins bonne que celle qu'elle remplace~: une instruction de journalisation
placée désormais dans la boucle critique, un conteneur qui alloue de la mémoire, un verrou pris sur un chemin qui était libre auparavant. Aucun de ces
changements ne modifie une sortie, donc aucun test fonctionnel n'échoue~; chacun ajoute pourtant des dizaines ou des centaines de nanosecondes au 99\up{e}~centile, c'est-à-dire
là où vit une stratégie sensible à la latence. Les détecter demande un autre type de test, qui mesure sur une machine qui
n'est jamais parfaitement calme, et décide si un écart est réel. Ce chapitre rassemble les tests dont un système à faible latence
a besoin, depuis ceux qui vérifient ce qu'il calcule (propriétés, différences, entrées aléatoires, scripts propres à la plateforme) jusqu'à celui
qui vérifie à quelle vitesse il calcule, et il s'achève sur la façon dont une version parvient en production~: progressivement, de manière réversible, et seulement sur
des machines qui ont été contrôlées.

\section{La pyramide des tests pour un système de trading}

\begin{definition}[Test par propriétés, test différentiel, test par fuzzing]\label{def:ll:testing-and-deploying-low-latency-systems:tests}
Un \emph{test par propriétés}\index{test par propriétés} génère de nombreuses entrées aléatoires et vérifie qu'une propriété énoncée est vraie
pour chacune d'elles, plutôt que de comparer une entrée à une sortie attendue unique. Un \emph{test
différentiel}\index{test différentiel} exécute deux implémentations d'une même spécification sur les mêmes entrées et vérifie que leurs
sorties sont identiques. Le \emph{test par fuzzing}\index{test par fuzzing} soumet à un programme un grand nombre d'entrées malformées ou mutées
afin de trouver celles qui le font planter, se bloquer ou violer ses invariants.
\end{definition}

Les composants de ce livre ont été testés de ces façons depuis le début. La passerelle du chapitre~21 est vérifiée par des propriétés sur
des entrelacements aléatoires face au moteur d'appariement (aucun ordre n'est exécuté au-delà de la quantité demandée, aucune position ne dépasse la limite)~; le
\omterm{def:ll:the-order-book-builder:builder}{constructeur de carnet} du chapitre~19 et les moteurs des chapitres~20 à 24 par des \omterm{def:ll:testing-and-deploying-low-latency-systems:tests}{tests différentiels}, puisque la référence en Python, les implémentations C++ et
Rust doivent concorder sur chaque jeu de données, octet pour octet ou empreinte pour empreinte. Le \omterm{def:ll:testing-and-deploying-low-latency-systems:tests}{test par propriétés} reprend l'idée de
QuickCheck (Claessen et Hughes, 2000)~: la propriété tient lieu de spécification, et le générateur explore des entrées que personne n'aurait
écrites à la main. La \Cref{fig:ll:testing-and-deploying-low-latency-systems:pyramid} range les types de tests dans l'ordre où une
version les rencontre~: beaucoup de tests peu coûteux en bas, exécutés à chaque modification, et quelques tests coûteux en haut, exécutés avant
la mise en production.

\begin{omfigure}
\begin{tikzpicture}[font=\scriptsize]
\foreach \y/\w/\t/\c in {0/11/{tests unitaires, de propriétés et différentiels~:\\ chaque modification, quelques secondes}/omDef,
  0.9/9.2/{fuzzing avec sanitizers~: chaque modification, quelques minutes}/omDef,
  1.8/7.4/{scénarios sur simulateur et\\ rejeux de journées enregistrées}/omMeth,
  2.7/5.6/{scripts de certification sur le\\ système de test de la plateforme}/omMeth,
  3.6/3.8/{barrière de performance~:\\ exécutions entrelacées}/omThm,
  4.5/2.2/{déploiement\\ progressif}/omThm} {
  \fill[\c!18,draw=\c!60] (-\w/2,\y) rectangle (\w/2,\y+0.8);
  \node[align=center] at (0,\y+0.4) {\t};
}
\draw[->] (6.3,0.2) -- node[right,align=left]{moins nombreux, plus lents,\\ plus proches de la production} (6.3,5.1);
\end{tikzpicture}

\omcaption{Les tests qu'un système de trading à faible latence rencontre avant la production. Les couches inférieures vérifient ce qu'il calcule et
s'exécutent à chaque modification~; les couches supérieures vérifient son comportement face à la plateforme et sa vitesse, et s'exécutent sur les candidats à
la mise en production.}\label{fig:ll:testing-and-deploying-low-latency-systems:pyramid}
\end{omfigure}

\section{Tests appuyés sur le simulateur de bourse}

Les tests unitaires s'arrêtent à la frontière du processus~; les bogues qui coûtent de l'argent se trouvent précisément à cette frontière~: un rapport qui arrive avant l'accusé de réception
de son ordre, une annulation qui croise une exécution, une session qui tombe au milieu d'une rafale, une plateforme qui rejette un message que la
firme jugeait valide. Le simulateur du livre~10 rend ces situations reproductibles~: la même graine donne le même entrelacement, une panne programmée
(une session coupée, une suspension, une pause) survient au même instant, et le moteur d'appariement applique les règles de la plateforme à la lettre.
Les \omterm{def:ll:testing-and-deploying-low-latency-systems:tests}{tests par propriétés} du chapitre~21 et la table de \omterm{def:ll:resilience:failover}{bascule} du chapitre~24 sont des tests appuyés sur le simulateur au sens où l'entend ce texte~: un scénario, un ensemble de
graines ou de points de coupure, et des propriétés vérifiées après chaque exécution. Les rejeux de journées enregistrées (chapitre~23) les complètent~: ils vérifient
qu'une nouvelle construction prend les décisions que prenait l'ancienne, sauf là où une modification était voulue.

Les entrées aléatoires atteignent des endroits que les scénarios ne couvrent pas. Le \Cref{lst:ll:testing-and-deploying-low-latency-systems:trusting} est un lecteur de somme de
contrôle affecté d'un bogue classique~: il fait confiance au champ BodyLength du message et lit la somme de contrôle à l'endroit indiqué par ce champ, sans
vérifier que le message est aussi long. Le fuzzer du chapitre mute les messages FIX de référence du chapitre~15 (un octet retourné,
un chiffre modifié, un champ dupliqué, une queue coupée) et recopie chacun dans un tampon de sa taille exacte, de sorte que
AddressSanitizer détecte toute lecture au-delà de celui-ci. Face à ce lecteur, il s'arrête au premier message tronqué, en 12 entrées au plus pour
chacune des dix graines. Face à l'analyseur propre à la firme, il a exécuté un million d'entrées sans lecture fautive. Il y a pourtant trouvé un vrai bogue,
par un invariant plutôt que par un plantage~: deux mutations, un tag 11 transformé en 10 et un numéro de séquence augmenté de un, ont laissé la
somme de contrôle inchangée, et l'analyseur a accepté un message dont un champ CheckSum se trouvait au milieu de son corps. Les analyseurs Python, C++ et Rust
de \texttt{firm.fixengine} le refusent désormais, et un test conserve ce message.

\omcode{code/low-latency/25-testing-and-deploying-low-latency-systems/cpp/ll_fuzz.hpp}{67}{78}{Un lecteur de somme de contrôle qui fait confiance à
BodyLength~: le premier message tronqué trouvé par le fuzzer le fait lire au-delà du tampon.}\label{lst:ll:testing-and-deploying-low-latency-systems:trusting}

\section{Certification par la bourse et conformité}

\begin{definition}[Certification par la bourse]\label{def:ll:testing-and-deploying-low-latency-systems:cert}
La \emph{certification par la bourse}\index{certification par la bourse} est la procédure par laquelle une plateforme vérifie, avant d'accorder l'accès à
la production, que le système d'un client traite correctement ses protocoles~: le client exécute un script de scénarios sur le système de test
de la plateforme, qui contrôle chaque message et chaque réponse.
\end{definition}

Les plateformes certifient les systèmes qui se connectent à elles, et les régulateurs demandent aux firmes de tester leur conformité à chaque plateforme avant
de négocier et après chaque changement important (l'encadré daté). Un script de certification est une liste d'étapes, chacune étant un message et les
réponses qu'il doit produire~: un ordre accepté, un remplacement, une annulation, une annulation arrivée trop tard, un identifiant réutilisé, un prix
hors du pas de cotation, un ordre exécuter-ou-annuler sans rien à exécuter, un ordre post-only qui croiserait le carnet, une
\omterm{def:ll:pre-trade-risk-and-kill-switches:mass}{annulation de masse} (mass cancel), une déconnexion qui annule les ordres au repos. Le lanceur de \texttt{firm.perfgate} joue un tel script contre le moteur d'appariement du livre~10
(\cref{lst:ll:testing-and-deploying-low-latency-systems:certify})~; les dix-sept étapes de la construction réussissent, et un script qui attend
le mauvais comportement d'une seule étape échoue à cette étape et à aucune autre. Le même lanceur, pointé vers une nouvelle version du
simulateur de bourse ou de la passerelle de la firme, constitue un test de conformité au sens du régulateur.

\omcode{code/firm/perfgate/firm_perfgate.py}{121}{142}{Le lanceur de certification~: chaque étape envoie un message et compare les
réponses, dans l'ordre, avec celles du script.}\label{lst:ll:testing-and-deploying-low-latency-systems:certify}

\begin{dated}{2026-09}{Test de conformité et certification}\label{dat:ll:testing-and-deploying-low-latency-systems:cert}
Dans l'Union européenne (consulté en septembre 2026), le Commission Delegated Regulation (EU) 2017/589 exige des firmes qu'elles établissent
des méthodologies pour développer et tester les systèmes de trading algorithmique avant leur déploiement ou une mise à jour substantielle~; qu'elles testent
leur conformité au système d'une plateforme de négociation lorsqu'elles deviennent membres, lorsqu'elles se connectent par accès sponsorisé pour la première fois,
lorsque les systèmes de la plateforme changent de manière importante, et avant de déployer ou de mettre à jour de manière importante un algorithme, en vérifiant
qu'il interagit avec la logique d'appariement de la plateforme comme prévu~; et qu'elles testent dans un environnement séparé de la production. CME Group
exige que les systèmes des clients qui envoient des ordres à sa plateforme Globex, ou qui traitent les données de marché qui en proviennent, soient certifiés avec AutoCert+,
son outil de test automatisé.
\end{dated}

\section{Barrières de régression de performance}

\begin{definition}[Barrière de régression de performance]\label{def:ll:testing-and-deploying-low-latency-systems:gate}
Une \emph{barrière de régression de performance}\index{barrière de régression de performance} est un test qui exécute un indice de référence sur la construction en cours
et sur un candidat, sur la même machine, et qui refuse le candidat si une statistique de latence se dégrade de plus que la
tolérance énoncée, avec un niveau de confiance énoncé.
\end{definition}

La \Cref{fig:ll:testing-and-deploying-low-latency-systems:runs} montre 800
essais du benchmark de la porte (le moteur de stratégie sur la journée à petit tick du chapitre~19, mesuré par événement), en alternant entre le
build courant et un build porteur d'une régression plantée. Le 99\up{e}~percentile d'un essai varie de dizaines de pour cent d'un essai à
l'autre sous l'effet de l'activité propre au système d'exploitation et à l'hyperviseur~; la régression, qui ajoute environ 18~\% à la médiane des
essais, est invisible dans
n'importe quelle paire isolée. La porte est conçue pour résister à ce bruit. Elle entrelace les essais (A, B, B, A, \dots), de sorte qu'une lente dérive de la
machine touche de la même façon les deux builds~; elle retient une statistique par essai, le 99\up{e}~percentile, et compare leurs médianes sur l'ensemble des essais~;
elle décide par un test de permutation, qui ne suppose rien sur la distribution, et rapporte un intervalle de bootstrap de la hausse
(\cref{lst:ll:testing-and-deploying-low-latency-systems:compare})~; et elle ne signale une régression que si la hausse est à la fois significative
et supérieure à une tolérance de 2~\%, afin qu'un changement réel mais négligeable ne bloque pas une mise en production.

\omcode{code/firm/perfgate/firm_perfgate.py}{62}{83}{La décision de la porte~: la hausse de la médiane des 99\up{e}~percentiles, un test de permutation unilatéral,
un intervalle de bootstrap et une tolérance.}\label{lst:ll:testing-and-deploying-low-latency-systems:compare}

\begin{omfigure}
\begin{tikzpicture}
\begin{axis}[width=11cm,height=5.4cm,xmin=0,xmax=800,ymin=150,ymax=700,xlabel={essai (entrelacé~: A B B A \dots)},
  ylabel={99\up{e}~percentile de l'essai (ns)},tick label style={font=\small},label style={font=\small},grid=major,
  grid style={gray!25},legend pos=north east,legend style={font=\scriptsize},table/col sep=comma]
\addplot[only marks,mark=*,mark size=0.7pt,omDef,restrict expr to domain={\thisrow{isb}}{0:0}]
  table[x=run,y=p99_ns]{figdata/low-latency/25-testing-and-deploying-low-latency-systems/gate_runs_plot.csv};
\addlegendentry{A~: build courant}
\addplot[only marks,mark=*,mark size=0.7pt,omThm,restrict expr to domain={\thisrow{isb}}{1:1}]
  table[x=run,y=p99_ns]{figdata/low-latency/25-testing-and-deploying-low-latency-systems/gate_runs_plot.csv};
\addlegendentry{B~: régression plantée}
\end{axis}
\end{tikzpicture}

\omcaption{Le 99\up{e}~percentile de chacun des 800 essais du benchmark de la porte (18\,000 événements mesurés par essai), en alternant entre
le build courant et un build porteur d'une régression plantée, sur un portable (Intel Core Ultra 7 155H, WSL2), la machine étant par ailleurs inactive.
Données~: \texttt{bench\_gate.py}.}\label{fig:ll:testing-and-deploying-low-latency-systems:runs}
\end{omfigure}

Combien d'essais la porte exige-t-elle~? La \Cref{fig:ll:testing-and-deploying-low-latency-systems:power} répond à partir du bruit propre
à la machine. Elle trace des portes de $n$ essais par côté, tirés des échantillons mesurés~: deux tirages disjoints dans les essais du build courant donnent
le taux de fausses alarmes, qui reste voisin du niveau du test, soit 5~\% (de 3 à 9~\% sur une centaine de portes)~; deux tirages dont l'un est décalé de
5~\% de la médiane donnent le taux de détection d'une hausse de 5~\%~; et la régression plantée, comparée au build courant, donne son
taux de détection. La régression plantée, d'environ 18~\%, est détectée dans 99~\% des cas avec 40 essais par côté. Une hausse de 5~\% en exige bien
davantage~: 73~\% à 120 essais, 84~\% à 160 et 91~\% à 200, soit environ 200 essais par côté pour 90~\%~; l'approximation normale de
\texttt{runs\_needed}, à partir de la dispersion des essais d'environ 14~\%, en donne 220. À environ 30 millisecondes par essai, cela représente une douzaine
de secondes de benchmark par porte sur ce portable, et sur
une machine proche de la production, dotée de cœurs isolés (chapitre~13), dont les essais varient moins, il en faut bien moins.

\begin{omfigure}
\begin{tikzpicture}
\begin{axis}[width=8.5cm,height=5.6cm,xmode=log,xmin=4,xmax=250,ymin=0,ymax=1.05,
  xlabel={essais par côté},ylabel={fraction des portes qui signalent B},tick label style={font=\small},label style={font=\small},
  grid=major,grid style={gray!25},legend pos=outer north east,legend style={font=\scriptsize,cells={anchor=west}},
  table/col sep=comma,xtick={5,10,20,40,80,200},xticklabels={5,10,20,40,80,200}]
\addplot[omThm,thick,mark=*] table[x=runs,y=detect_planted]{figdata/low-latency/25-testing-and-deploying-low-latency-systems/gate.csv};
\addlegendentry{régression plantée (d'environ 18~\%)}
\addplot[omMeth,thick,mark=square*] table[x=runs,y=detect_5pct]{figdata/low-latency/25-testing-and-deploying-low-latency-systems/gate.csv};
\addlegendentry{une hausse de 5~\%}
\addplot[omDef,thick,mark=triangle*] table[x=runs,y=false_alarm]{figdata/low-latency/25-testing-and-deploying-low-latency-systems/gate.csv};
\addlegendentry{aucune variation (fausses alarmes)}
\addplot[gray,dashed] coordinates {(4,0.9) (250,0.9)};
\end{axis}
\end{tikzpicture}

\omcaption{Les taux de détection et de fausses alarmes de la porte en fonction du nombre d'essais par côté, à partir de 100 portes tirées des
essais mesurés de la \cref{fig:ll:testing-and-deploying-low-latency-systems:runs} pour chaque point~; la ligne en pointillés correspond à une puissance de 90~\%.
Données~: \texttt{fig\_gate.py} (déterministe, à partir des échantillons mesurés).}\label{fig:ll:testing-and-deploying-low-latency-systems:power}
\end{omfigure}

\section{Déploiements progressifs et retour arrière}

Une version qui a franchi la porte subit encore une mise en production graduelle~: le déploiement canari et le déploiement progressif du
livre~7 (chapitre~21), avec du trading simulé lorsque la stratégie le permet. Les systèmes à faible latence ajoutent trois réflexes. L'hôte est contrôlé
avant le démarrage du processus~: \texttt{firm.perfgate.preflight} exécute l'audit de réglage du chapitre~13 contre le plan de placement des cœurs
et interrompt le déploiement à la moindre anomalie, car une stratégie testée sur des cœurs isolés et déployée sur un hôte dont les interruptions tombent sur
ces cœurs n'est pas la stratégie qui a été testée (ce portable échoue sur 7 des 11 règles de l'audit, raison pour laquelle ses mesures portent
la mention «~pas de cœurs isolés~»). Le nouveau build tourne d'abord à côté de l'ancien, sur les mêmes entrées, en mode shadow~: le rejeu
du chapitre~23 rendu réel, dont les sorties sont comparées puis écartées. Et le retour arrière est un déploiement comme les autres, préparé et
testé avant d'être nécessaire, car le moment où il le devient est celui où l'équipe dispose du moins de temps~; la \omterm{def:ll:resilience:failover}{bascule} (failover)
du chapitre~24 illustre la même leçon pour les machines.

\section{Tutoriel~: fuzzer, certifier, encadrer}\label{tut:ll:testing-and-deploying-low-latency-systems:tut}

\begin{tutorial}
\textbf{Objectif.} Fuzzer l'analyseur FIX sous les sanitiseurs, lancer un script de certification contre le simulateur et soumettre à la porte une
régression plantée sur ce portable. \textbf{État final~:} les résultats du fuzzing, une certification réussie et la
\cref{fig:ll:testing-and-deploying-low-latency-systems:power}.
\begin{tutsteps} \sloppy
\item \textbf{Fuzzing.} \texttt{ll\_fuzz\_test.cpp} exécute 200\,000 mutations à travers l'analyseur de la maison et vérifie que chaque
message accepté tient dans son tampon et porte une somme de contrôle correcte~; \texttt{python bench\_fuzz.py} compile le fuzzer avec AddressSanitizer et
UndefinedBehaviorSanitizer et le lance contre les deux analyseurs.
\item \textbf{Certification.}\newline \texttt{firm\_perfgate.certify(load\_script(\textquotedbl{}data/cert\_script.json\textquotedbl{}))}
joue les dix-sept étapes contre le moteur du livre~10.
\item \textbf{Porte.} \texttt{python bench\_gate.py} mesure 400 paires entrelacées~; \texttt{python fig\_gate.py} les transforme
en taux de détection et de fausses alarmes.
\item \textbf{Contrôle préalable.} \texttt{firm\_perfgate.preflight} sur les deux hôtes de test de l'audit de réglage~: l'hôte réglé réussit, l'hôte
mal réglé échoue.
\end{tutsteps}
\textbf{Pour aller plus loin.} Faire reposer la porte sur la médiane des médianes des essais au lieu de leurs 99\up{e}~percentiles et comparer le nombre
d'essais nécessaires~; ajouter une étape au script de certification pour le bridage de la plateforme.
\end{tutorial}

\section{Construction~: la porte de performance}\label{bld:ll:testing-and-deploying-low-latency-systems:perfgate}

\begin{build}
\bfield{Rôle} Les dernières vérifications avant la production pour chaque composant du livre~: sa vitesse par rapport au build courant,
son comportement face au script de la plateforme, et l'hôte sur lequel il tournera.
\bfield{Interface} Python \texttt{firm\_perfgate}~: \texttt{interleave}, \texttt{run\_pairs(cmd\_a, cmd\_b, n\_pairs, parse)},
\texttt{compare(a, b, tolerance, alpha)} renvoyant un \texttt{Verdict} avec \texttt{rise}, \texttt{p\_value}, \texttt{low},
\texttt{high}, \texttt{regression} et \texttt{report()}~; \texttt{power}, \texttt{runs\_needed}~; \texttt{certify(script)},
\texttt{load\_script}~; \texttt{preflight(root, plan)}.
\bfield{Règles} Essais entrelacés, jamais tous les A puis tous les B~; une statistique par essai~; un test de permutation et une tolérance~; le
nombre d'essais fixé à partir du bruit mesuré de la machine~; aucun déploiement sur un hôte qui échoue à l'audit.
\bfield{Tests d'acceptation} \texttt{code/firm/perfgate/}~: une hausse nette doit être signalée, alors que ni le bruit ni une hausse sous la tolérance ne doivent l'être~; des fausses
alarmes proches du niveau du test sur du bruit synthétique~; une puissance qui croît avec le nombre d'essais~; le script de certification réussi,
et échouant précisément à l'étape altérée~; le contrôle préalable réussi sur l'hôte de test réglé et échouant sur l'hôte mal réglé.
\bfield{Pour aller plus loin} Une porte sur toute la distribution de latence (un test à deux échantillons sur des échantillons regroupés, avec attention à la dépendance au sein des
essais)~; des étapes de certification pour les sessions et les bridages~; la porte en intégration continue sur une machine dédiée.
\end{build}

\begin{omsources}
\item K.~Claessen et J.~Hughes, «~QuickCheck: a lightweight tool for random testing of Haskell programs~», \emph{ICFP}, 2000.
\item Commission Delegated Regulation (EU) 2017/589 (RTS 6), articles 5 à 7.
\item CME Group Client Systems Wiki, \emph{Client Application Testing and Certification}.
\end{omsources}

\section{Exercices}

\begin{exercise}[$\star$]\label{exo:ll:testing-and-deploying-low-latency-systems:1}
Énoncer deux propriétés qu'un test sur des entrelacements aléatoires doit vérifier pour la \omterm{def:ll:order-gateway-and-order-management:gateway}{passerelle d'ordres} du chapitre~21, et une
qu'il ne peut pas vérifier sans le \omterm{def:ll:order-gateway-and-order-management:dropcopy}{drop copy} de la plateforme.
\end{exercise}

\begin{exercise}[$\star$]\label{exo:ll:testing-and-deploying-low-latency-systems:2}
Pourquoi le fuzzer copie-t-il chaque entrée dans un tampon de sa taille exacte avant de l'analyser~?
\end{exercise}

\begin{exercise}[$\star$]\label{exo:ll:testing-and-deploying-low-latency-systems:3}
Les exécutions de A et de B sont entrelacées A, B, B, A. Qu'est-ce qui ne va pas si toutes les exécutions de A viennent d'abord, puis toutes celles de B~?
\end{exercise}

\begin{exercise}[$\star\star$]\label{exo:ll:testing-and-deploying-low-latency-systems:4}
Avec des exécutions dont les 99\up{e}~centiles varient de 14~\%, combien d'exécutions par côté l'approximation normale demande-t-elle pour détecter une hausse de 10~\%
avec une puissance de 90~\%, puis une hausse de 2.5~\%~?
\end{exercise}

\begin{exercise}[$\star\star$]\label{exo:ll:testing-and-deploying-low-latency-systems:5}
Une barrière s'exécute 20 fois par jour à un niveau de 5~\%. Combien de fausses alarmes l'équipe doit-elle attendre sur un mois de 21 jours de bourse, et
que change la tolérance de 2~\%~?
\end{exercise}

\begin{exercise}[$\star\star$]\label{exo:ll:testing-and-deploying-low-latency-systems:6}
Deux mutations ont fait passer un tag de 11 à 10 et un numéro de séquence de 2 à 3, et la somme de contrôle n'a pas changé. Expliquer pourquoi, et
dire ce que la somme de contrôle FIX protège et ne protège pas.
\end{exercise}

\begin{exercise}[$\star\star\star$]\label{exo:ll:testing-and-deploying-low-latency-systems:7}
\emph{Programmation.} Ajouter des étapes de certification pour le bridage de la plateforme~: configurer le moteur avec un débit, envoyer une rafale et vérifier les
rejets de motif \texttt{T}. De quoi l'étape a-t-elle besoin que les autres n'exigeaient pas~?
\end{exercise}

\begin{exercise}[$\star\star\star$]\label{exo:ll:testing-and-deploying-low-latency-systems:8}
\emph{Trouver la faille.} «~Notre test de performance lance le benchmark une seule fois sur le candidat à la mise en production et échoue si le 99\up{e}~centile
dépasse de plus de 5~\% celui de la semaine dernière, que nous conservons dans un fichier.~»
\end{exercise}

\section{Problème~: quelques centaines de nanosecondes}

\begin{problem}[{Problème du week-end --- combien d'exécutions une barrière demande sur une machine bruitée}]\label{pb:ll:testing-and-deploying-low-latency-systems:1}
Utiliser les \cref{fig:ll:testing-and-deploying-low-latency-systems:runs,fig:ll:testing-and-deploying-low-latency-systems:power}
(\texttt{measured\_gate\_runs.csv}, \texttt{gate.csv}, \texttt{gate\_summary.csv}).

\medskip\noindent\textbf{Partie I --- Le bruit.}
\begin{enumerate}
\item Quelle est la médiane du 99\up{e}~centile des exécutions de la version courante, et de combien varient les exécutions (le coefficient de variation
robuste)~?
\item Qu'ajoute la régression plantée, à la médiane des exécutions~?
\item Pourquoi le 99\up{e}~centile d'une exécution est-il beaucoup plus bruité que sa médiane~?
\item Pourquoi une seule paire A/B est-elle inutile sur cette machine~?
\end{enumerate}

\medskip\noindent\textbf{Partie II --- La barrière.}
\begin{enumerate}[resume]
\item Pourquoi comparer des médianes entre exécutions plutôt que des moyennes~?
\item Que suppose le test de permutation, et que ne suppose-t-il pas~?
\item Quel est le taux de fausses alarmes pour chaque $n$, et pourquoi est-il proche du niveau du test~?
\item Combien d'exécutions par côté détectent la régression plantée dans 99~\% des cas~?
\end{enumerate}

\medskip\noindent\textbf{Partie III --- Puissance.}
\begin{enumerate}[resume]
\item Lire le taux de détection d'une hausse de 5~\% à 80, 120 et 160 exécutions par côté.
\item Interpoler le nombre d'exécutions pour une puissance de 90~\%.
\item Comparer avec l'approximation normale, et expliquer l'écart.
\item Combien de temps dure une barrière de cette taille sur cet ordinateur portable, à environ 30 millisecondes par exécution~?
\end{enumerate}

\medskip\noindent\textbf{Partie IV --- Le verdict.}
\begin{enumerate}[resume]
\item Énoncer le \emph{résultat nommé}~: le nombre d'exécutions entrelacées dont la barrière a besoin pour détecter une hausse de 5~\% du 99\up{e}~centile
avec une puissance de 90~\%, au bruit mesuré de cette machine.
\item Que changeraient les cœurs isolés (chapitre~13), et pourquoi~?
\item Pourquoi faut-il une tolérance en plus du test~?
\item De quoi aurait besoin une barrière sur la latence médiane, et pourquoi ne suffit-elle pas~?
\item Comment faire de la barrière une étape routinière de l'intégration continue sans que chaque modification attende~?
\item Quels types de régression aucun benchmark du seul moteur ne peut-il détecter~?
\item Pourquoi conserver les ensembles d'exécutions comme données versionnées~?
\item En une phrase~: qu'est-ce qui rend une barrière de performance fiable~?
\end{enumerate}
\end{problem}

\section{Questions d'entretien}

\begin{interviewq}[$\star$ \iqroles{développeur}]\label{iq:ll:testing-and-deploying-low-latency-systems:1}
Qu'est-ce qu'un \omterm{def:ll:testing-and-deploying-low-latency-systems:tests}{test par propriétés}~? Donnez un exemple pour un carnet d'ordres.
\end{interviewq}

\begin{interviewq}[$\star\star$ \iqroles{développeur}]\label{iq:ll:testing-and-deploying-low-latency-systems:2}
Comment mèneriez-vous un \omterm{def:ll:testing-and-deploying-low-latency-systems:tests}{test par fuzzing} sur un analyseur de données de marché, et que vérifieriez-vous en plus des plantages~?
\end{interviewq}

\begin{interviewq}[$\star\star$ \iqroles{développeur}]\label{iq:ll:testing-and-deploying-low-latency-systems:3}
Votre benchmark indique que la nouvelle version est plus lente de 3~\% au 99\up{e}~centile. Comment décidez-vous s'il faut la croire~?
\end{interviewq}

\begin{interviewq}[$\star\star$ \iqroles{développeur}]\label{iq:ll:testing-and-deploying-low-latency-systems:4}
Que teste la \omterm{def:ll:testing-and-deploying-low-latency-systems:cert}{certification par la bourse}, et que ne teste-t-elle pas~?
\end{interviewq}

\begin{interviewq}[$\star\star$ \iqroles{développeur, trader}]\label{iq:ll:testing-and-deploying-low-latency-systems:5}
Comment déploieriez-vous une nouvelle version d'un moteur de stratégie, et comment la feriez-vous revenir en arrière~?
\end{interviewq}

\begin{interviewq}[$\star\star\star$ \iqroles{développeur}]\label{iq:ll:testing-and-deploying-low-latency-systems:6}
Concevez le pipeline de tests et de mise en production d'un système de trading à faible latence, depuis un commit jusqu'à la production, et dites ce que chaque étape
détecte.
\end{interviewq}
